% !TeX root = ../../../main.tex

\chapter{Introducción a la Arquitectura lógica y física de la solución}
\label{cap:arquitectura-logica}

El capítulo 4 explica cómo la solución sostiene la operación distribuida de Puelche. La arquitectura lógica define las responsabilidades, los datos y los contratos que permiten ejecutar las capacidades del capítulo 3; la arquitectura física acredita su emplazamiento y continuidad, y la estrategia de centros de datos completa su recuperación. Esta descripción se relaciona con el modelo de datos del capítulo 5, el registro ADR y el Formulario T-11 de la arquitectura integrada.

\section{Arquitectura lógica}
\label{sec:arquitectura-logica}

Los anexos A--N detallan eventos, módulos, interfaces, reglas y correspondencias. O contiene las decisiones lógicas fechadas; P, las tecnologías y su actualización; Q y R, amenazas y controles; S, los puntos de vista y sus reglas de correspondencia; T, el desempeño; U, la evidencia documental; y V, las condiciones de aceptación que requieren intervención del CLIENTE. El archivo independiente de anexos comienza con un catálogo navegable: cada entrada identifica la sección que la utiliza y el requisito que respalda. Los identificadores se mantienen aunque cambie una página.

Este apartado desarrolla las responsabilidades y los contratos de la solución que deben corresponder al esquema y explicación de los apartados 3.3 y 3.4. La correspondencia se establece por capacidad: recepción e inventario (M1--M2), preventa y planificación (M3--M4), preparación y reparto (M5--M6), rendición y devoluciones (M7--M8), calidad y analítica (M9--M10), canal moderno y flota (M11--M12). El emplazamiento, las conexiones y su capacidad corresponden a 4.2; los centros de datos, a 4.3. Los catálogos extensos se entregan en los anexos 4.1-A a 4.1-N. El Anexo N identifica las capacidades del capítulo 3 disponible y la realización requerida de cada módulo; distingue esa cobertura lógica del cotejo completo que exige el capítulo integrado.

En Puelche, un pedido debe poder tomarse en una ruta sin cobertura, prepararse durante la noche y llegar al cliente con su lote y su evidencia de entrega identificados. La arquitectura parte de esa continuidad: cada operación se registra donde ocurre y se reconcilia cuando vuelve la conexión. Para ordenar las responsabilidades, la solución se organiza en las ocho capas exigidas por el numeral 2.1 de las Bases Técnicas Transversales (Pontificia Universidad Católica de Valparaíso [PUCV], 2026b, cap.~2, p.~6; RT-02.01).

